% =========================================================
\subsection{Falla detectada en la primera secuencia con instrumento:
falsos positivos del testigo de adquisición}
\label{sec:falla_testigo}
% =========================================================

La validación del testigo de adquisición nueva descrita en la
\cref{sec:verificacion_adquisicion} se hizo con capturas manuales; su
comportamiento en los modos automáticos solo se había ejercitado sin
hardware. Las sesiones del 3 de noviembre de 2026, adquiridas con un
TDS5054B (número de serie B022263) y el láser emitiendo, fueron la
primera ocasión de observarlo en secuencias reales. El resultado fue
una falla del propio mecanismo de detección: marcó como inválidas
capturas que no lo eran.

\subsubsection*{Lo que se observó}

En la sesión \texttt{20261103\_194433}, en modo tiempo con 100
promedios, las siete capturas consecutivas \texttt{m0003} a
\texttt{m0009} quedaron marcadas como \texttt{osciloscopio sin
adquisición nueva}, aunque cada una registra en
\texttt{adquisiciones\_promediadas} el promedio completo de 100. En la
sesión \texttt{20261103\_200010}, en modo temperatura, las capturas
\texttt{m0006} y \texttt{m0007} recibieron la misma marca con 444 y
416 adquisiciones en ventanas de \SI{44.5}{\second} y
\SI{42.1}{\second}. Esa tasa, de 10.0 y 9.9 adquisiciones por segundo,
es la misma que la de las capturas aceptadas sin marca en esa
secuencia, entre 9.9 y 10.0, y corresponde a la frecuencia de
repetición del láser. Lo mismo ocurrió con la captura \texttt{m0006}
de la sesión \texttt{20261103\_190015}, marcada con 532
adquisiciones en \SI{53.6}{\second}.

Dos comprobaciones independientes del contador confirman que se trata
de falsos positivos. Primera: el registro de cada captura marcada es
distinto del de la anterior. A diferencia de lo observado el 22 de
septiembre, donde esa diferencia no probaba nada porque se debía a un
cambio de escala con el instrumento detenido, aquí las escalas vertical
y horizontal de cada captura marcada son las mismas que las de la
captura anterior. Segunda: en
las mismas sesiones el testigo sí marcó las dos capturas que en
efecto carecían de adquisición válida. En la sesión
\texttt{20261103\_190015}, la captura \texttt{m0008} registra cero
adquisiciones en una ventana de \SI{69.6}{\second} y su registro es
idéntico, bit a bit, al de \texttt{m0007}: el láser había dejado de
emitir. En la sesión \texttt{20261103\_194433}, la captura
\texttt{m0010} se marcó con \texttt{promedio incompleto (84 de 100)}.

\subsubsection*{Diagnóstico}

Todas las capturas marcadas por error comparten un rasgo: su conteo
final es igual o menor que el de la captura anterior. El testigo se
diseñó sobre una propiedad medida el 22 de septiembre: que
\texttt{ACQ:STATE RUN} no reinicia el contador de adquisiciones. Esa
medición se hizo con el osciloscopio adquiriendo de forma continua
(\texttt{ACQ:STOPAFTER RUNSTOP}), donde la orden llega a un
instrumento que ya está en marcha. En los modos automáticos, en
cambio, el osciloscopio opera con \texttt{ACQ:STOPAFTER SEQUENCE} y
queda detenido entre capturas, de modo que cada \texttt{ACQ:STATE RUN}
inicia una secuencia nueva y reinicia el contador. El reinicio, sin
embargo, no es inmediato: la lectura que el testigo toma justo después
de la orden devuelve todavía el conteo de la captura anterior. En
estos modos el testigo solo recibe esa lectura inicial y la final, por
lo que interpreta el reinicio como una caída sin subida posterior y
declara la captura sin adquisición nueva. En modo tiempo el efecto es
sistemático, porque toda captura termina exactamente en 100. En modo
temperatura es intermitente: en las capturas \texttt{m0004} y
\texttt{m0011} de la sesión \texttt{20261103\_200010}, cuyo conteo
también fue menor que el anterior, el reinicio ya era visible en la
primera lectura y no hubo marca, lo que indica que el retraso del
reinicio varía.

\subsubsection*{El error en sentido contrario}

El análisis de las mismas sesiones mostró que el testigo tampoco
detecta el caso opuesto. En la sesión \texttt{20261103\_200010}, las
capturas \texttt{m0008} a \texttt{m0013} se aceptaron sin marca con
tasas de entre 13 y 69 adquisiciones por segundo, superiores a la
frecuencia del láser: el promedio de esas capturas incluye
adquisiciones ajenas a su ventana de temperatura. La causa de ese
exceso no quedó establecida con los datos disponibles, porque el
registro de sesión no conserva la configuración del disparo del
osciloscopio. Lo que sí quedó establecido es la carencia del testigo:
exige que el contador suba, pero no compara cuánto subió con lo que la
duración de la ventana permite. Por la misma razón, en la sesión
\texttt{20261103\_190015} se aceptaron sin marca dos capturas en cuyas
ventanas cambió la escala vertical: \texttt{m0003}, que reunió solo 44
adquisiciones en \SI{45}{\second}, y \texttt{m0007}. En modo
temperatura el testigo no comparaba escalas, a diferencia de la
captura manual.

\subsubsection*{Nivel de detección}

La falla se manifestó en ejecución, pero como una alarma falsa: el
sistema avisó, y lo que avisaba era incorrecto. Su diagnóstico, y el
hallazgo del error contrario, solo fueron posibles en la verificación
posterior, al confrontar el contador registrado con la duración de
cada ventana y con el contenido de los registros. El caso ilustra un
riesgo propio de la automatización de instrumentos: una propiedad
medida en el equipo es válida para el régimen en que se midió, y el
testigo la extendió a un régimen distinto sin volver a medirla.
